\documentclass[10pt,a4paper]{amsart}
\usepackage[margin=19mm]{geometry}
\usepackage{amsmath,amssymb,mathtools}
\usepackage[scr=rsfs,cal=euler]{mathalfa}
\usepackage{xfrac}
\usepackage[unicode,hidelinks,bookmarks=false]{hyperref}
\newcommand{\ie}{i.e.\ }
\newcommand{\cf}{cf.\ }
\newcommand{\resp}{respectively\ }
\newcommand{\Subsection}{\subsection}
\providecommand{\nobreakdash}{\nobreak\hbox{-}}
\setlength{\emergencystretch}{2em}
\multlinegap0pt
\usepackage{bm}
\usepackage{bbm}
\usepackage{dsfont}
\usepackage{xspace}
\usepackage{cancel}
\usepackage{mathtools}
\usepackage{amsthm}
\makeatletter
\@ifundefined{theorem}{\newtheorem{theorem}{Theorem}}{}
\@ifundefined{lemma}{\newtheorem{lemma}[theorem]{Lemma}}{}
\@ifundefined{proposition}{\newtheorem{proposition}[theorem]{Proposition}}{}
\makeatother
\DeclareMathOperator{\rad}{rad}
\title[Must a primitive non-deficient number have a component not m]{Must a primitive non-deficient number have a component not much larger than its radical?}
\author{Joshua Zelinsky}
\date{Source: arXiv:2005.12115v5 (CC BY 4.0)}
\begin{document}
\maketitle

\noindent\emph{Source and licence.} Must a primitive non-deficient number have a component not much larger than its radical?. Version arXiv:2005.12115v5 (CC BY 4.0), distributed under the Creative Commons Attribution 4.0 International licence, as recorded in the arXiv licence metadata for this exact version and shown on the abstract page https://arxiv.org/abs/2005.12115v5. Full rights notice: licenses/arxiv-2005.12115-CC-BY-4.0.txt.\par\smallskip

\noindent\textit{Editorial scope.} Counted unit: the Proposition whose source label is \texttt{pklessthansquarerootofR} in the article (source0.tex lines 328--330, source label \texttt{pklessthansquarerootofR}), delivered together with the article's own complete original proof of it (source0.tex lines 331--373, one \texttt{proof} environment of 43 lines). No mathematics is written, repaired or re-derived: statement and proof are the article's own text line for line. Required context retained uncounted: GeneralizedPuchtalemma (lines 224--226). Every definition, display and label of those lines is retained unchanged. Omitted from this package and not redistributed: 1--223, 227--327, 374--463. Nothing inside a retained excerpt is omitted or altered: every retained body is delivered exactly as the article's lines are written, so no omission receipt of any kind accompanies this package. Citations: the retained excerpts carry no citation command at all, so no bibliography entry of the article is transcribed and no reference is redistributed. Reference adaptation: none was needed and none was made; every reference inside the delivered unit resolves inside it, so no unresolved reference or citation is produced. Printed slips kept literal: no printed word, symbol, hypothesis or number of the counted statement or of its proof was altered. Layout and class: the article's own class is \texttt{article} with packages including color, amssymb, amsmath, amsthm, amsfonts, amscd, graphicx, xcolor, float, hyperref, fullpage, graphics, latexsym, epsf; the wrapper is the lane's pinned \texttt{amsart} editorial wrapper at 10pt on A4, which restates the theorem-like environments the retained text needs and adds the article's own macro definitions that the retained text needs (\texttt{\textbackslash rad}), with extra wrapper packages (bm, bbm, dsfont, xspace, cancel, mathtools). The delivered numbering is the unit's own. Assets: no retained span contains a figure environment, a graphics-inclusion command, a PDF-page inclusion, a table or a TikZ picture; no image file is copied into this package, the package declares no asset dependency and no image file is redistributed with it. Licence: version arXiv:2005.12115v5 of this article is distributed under the Creative Commons Attribution 4.0 International licence, as recorded in the arXiv licence metadata for this exact version and shown on the abstract page https://arxiv.org/abs/2005.12115v5. A full rights notice is delivered at licenses/arxiv-2005.12115-CC-BY-4.0.txt. Characters: every retained excerpt is pure ASCII, so the delivered engine, XeLaTeX with its default Latin Modern text font, reports no missing character.
\begin{lemma}\label{GeneralizedPuchtalemma} Let $n$ be a primitive non-deficient number with $n=p_1^{a_1}p_2^{a_2} \cdots p_k^{a_k}$ with $p_1 < p_2 < \cdots < p_k$. Assume further that $(p_1,k) \neq (3,3)$. Let $R=\rad(n) = p_1p_2\cdots p_k$. Assume further that $$H(n) \geq 2 + \alpha$$ for some positive $\alpha <1$. Then there exists an $i$, $1 \leq i \leq k$ such that 
$$p_i^{a_i+1} < \max\left(\frac{2(k+2+p_1)}{\alpha },k(k+1)\right).$$
\end{lemma}

\begin{proposition}  Let $n$ be an odd  primitive non-deficient number with $n=p_1^{a_1}p_2^{a_2} \cdots p_k^{a_k}$ with $p_1 < p_2 < \cdots < p_k$.
Either \begin{equation}p_k < \sqrt{R}\label{pklessthansquarerootofR}\end{equation} or there exists an $i$ such that \begin{equation}p_i^{a_i+1} < (k+2+p_1)(p_k-1).\label{mincomponentnotmuchbiggerthanlargestprimefactor}\end{equation}\label{dichotomywithpk}
\end{proposition}

\begin{proof} 

If $p_1=3$ and $k=3$, it is easy to check the short finite list that all of these satisfy at least one of the two bounds. We thus may assume that $(p_1, k) \neq (3,3)$.

As before, if $n$ is non-deficient we must have
$$H(n) = \prod_{i=1}^k \frac{p_i}{p_i-1} >2, $$
and thus 
$$\frac{p_1p_2p_3 \cdots p_k}{(p_1 -1)(p_2-1) \cdots (p_k-1)} -2 >0 $$
which means that \begin{equation}p_1p_2 \cdots p_k  = 2(p_1-1)(p_2-1)\cdots (p_k-1) +x \label{xintroduction} \end{equation} for some $x$. Note that $x$ is itself a positive integer. 
Now,  Equation \eqref{xintroduction} modulo $p_k$ becomes
$$0 \equiv 2(p_1-1)(p_2-1)\cdots(p_{k-1}-1)(-1) + x \pmod{p_k}.$$
So $p_k \mid -2(p_1-1)(p_2-1)\cdots(p_{k-1}-1) + x$.
If $-2(p_1-1)(p_2-1)\cdots(p_{k-1}-1) + x$ is negative, then
$p_k \mid 2(p_1-1)(p_2-1)\cdots(p_{k-1}-1) -x. $
Thus,
\begin{equation}p_k \leq 2(p_1-1)(p_2-1)\cdots(p_{k-1}-1) -x < 2(p_1-1)(p_2-1)\cdots(p_{k-1}-1) \label{initialpkinequalitytogetRform} \end{equation}
We have then 
\begin{equation} 2(p_1-1)(p_2-1)\cdots(p_{k-1}-1)= 2R\frac{p_1-1}{p_1}\frac{p_2-1}{p_2} \cdots \frac{p_{k-1}-1}{p_{k-1}}\frac{1}{p_k} \label{intermediateboundtogetRform}
\end{equation}

Since $n$ is non-deficient we have
$$\frac{p_1-1}{p_1}\frac{p_2-1}{p_2} \cdots \frac{p_{k-1}-1}{p_{k-1}}\frac{p_{k-1}}{p_{k}} <2, $$

and so 
\begin{equation} \frac{p_1-1}{p_1}\frac{p_2-1}{p_2} \cdots \frac{p_{k-1}-1}{p_{k-1}} < \frac{1}{2}\frac{p_k}{p_k-1} \label{withoutpkproductclosetoonehalf}.
\end{equation}

We then combine inequalities \eqref{initialpkinequalitytogetRform}, \eqref{intermediateboundtogetRform} and \eqref{withoutpkproductclosetoonehalf} to get that

$$p_k \leq \frac{2R}{p_k}\frac{1}{2}\frac{p_k}{p_k-1}$$ which implies that
\begin{equation}p_k (p_k-1) < R.\label{pktimespkminus1lessthanR}\end{equation}

Since $R$ is a multiple of $p_k$, we must have $$p_k(p_k-1) \leq R-p_k$$ which implies that $p_k < \sqrt{R}$.

Now, assume that $$p_k \geq  \sqrt{R}.$$ Thus, by our above reasoning we must that $-2(p_1-1)(p_2-1)\cdots(p_{k-1}-1) + x$  is non-negative. Then we have that $x \geq 2(p_1-1)(p_2-1)\cdots(p_{k-1}-1)$.

Thus, we have that
\begin{equation}p_1p_2 \cdots p_k  \geq  2(p_1-1)(p_2-1)\cdots (p_k-1) + 2(p_1-1)(p_2-1)\cdots(p_{k-1}-1) \label{Wegettousex},\end{equation}
and so
$$\prod_{i=1}^k \frac{p_i}{p_i-1} \geq 2 +\frac{2(p_1-1)(p_2-1)\cdots(p_{k-1}-1)}{(p_1-1)(p_2-1)\cdots(p_k-1)} =  2 + \frac{2}{p_k-1}.   $$

The rest of the proof then continues from invoking Lemma \ref{GeneralizedPuchtalemma} and proceeding as before. 
\end{proof}

\end{document}
